\documentclass[11pt]{article}

\usepackage{amsmath}
\usepackage{amssymb}
\usepackage{amsthm}
\usepackage{geometry}
\usepackage{mathrsfs}
\usepackage{mathtools}
\usepackage{tikz}
\usepackage[colorlinks=true, linkcolor=blue, bookmarks=true]{hyperref}

\geometry{letterpaper, margin=1in}

\renewcommand{\qedsymbol}{$\blacksquare$}

\title{\textbf{Advanced Mathematical Logic and Set Theory}\\ \Large Rigorous Practice Problem Set}
\author{Elite Professor of Pure Mathematics}
\date{\today}

\begin{document}

\maketitle

\tableofcontents
\newpage

\section{Concept Checklist \& Question Mapping}
\label{sec:mapping_start}
This problem set is strictly derived from the theoretical frameworks, logical structures, and abstract modeling paradigms detailed in the provided text. Every problem maps directly to a core competency required for mastery of set-theoretic proofs.

\begin{itemize}
    \item \textbf{C1: Set Equality via Double Inclusion} (Problems 1--4)
    \item \textbf{C2: Biconditional Set Properties} (Problems 5--8)
    \item \textbf{C3: Empty Set, Disjointness, and Contrapositive/Contradiction} (Problems 9--12)
    \item \textbf{C4: Sets Defined by Number Theory (Modular Arithmetic \& Divisibility)} (Problems 13--16)
    \item \textbf{C5: Power Sets ($\mathcal{P}(A)$) and Element vs. Subset Logic} (Problems 17--20)
    \item \textbf{C6: Generalized Set Operations (Infinite Unions/Intersections)} (Problems 21--22)
    \item \textbf{C7: Abstract Modeling: Machine Learning (Symmetric Difference)} (Problems 23--26)
    \item \textbf{C8: Abstract Modeling: Quantitative Finance (Probability Spaces \& Hitting Times)} (Problems 27--28)
    \item \textbf{C9: Diagnostic Testing: Identifying Logical Flaws} (Problems 29--30)
\end{itemize}

\newpage

\section{Problem Sets}
\label{sec:problems}

\subsection{Set Equality via Double Inclusion}
\noindent \textbf{Problem 1.} Let $A$ and $B$ be sets. Prove that $(A \cup B) - (A \cap B) = (A - B) \cup (B - A)$. \\
\textbf{Problem 2.} Let $A, B,$ and $C$ be sets. Prove that $A - (B - C) = (A - B) \cup (A \cap C)$. \\
\textbf{Problem 3.} Let $A, B,$ and $C$ be sets. Prove that $(A - B) - C = A - (B \cup C)$. \\
\textbf{Problem 4.} Let $A, B,$ and $C$ be sets. Prove that $A \cap (B - C) = (A \cap B) - (A \cap C)$.

\subsection{Biconditional Set Properties}
\noindent \textbf{Problem 5.} Let $A$ and $B$ be sets. Prove that $A \cup B = B$ if and only if $A \subseteq B$. \\
\textbf{Problem 6.} Let $A$ and $B$ be sets. Prove that $A - B = \emptyset$ if and only if $A \subseteq B$. \\
\textbf{Problem 7.} Let $A$ and $B$ be sets. Prove that $A \cap B = A \cup B$ if and only if $A = B$. \\
\textbf{Problem 8.} Let $A$ and $B$ be subsets of a universal set $U$. Prove that $A \subseteq B$ if and only if $\overline{B} \subseteq \overline{A}$.

\subsection{Empty Set and Disjointness Proofs}
\noindent \textbf{Problem 9.} Let $A, B,$ and $C$ be sets. Prove that if $A \cap B \neq \emptyset$ and $A \subseteq C$, then $C \cap B \neq \emptyset$. \\
\textbf{Problem 10.} Prove by contradiction: For any sets $A$ and $B$, if $A \cup B = \emptyset$, then $A = \emptyset$ and $B = \emptyset$. \\
\textbf{Problem 11.} Let $A$ and $B$ be sets. Prove that $A$ and $B - A$ are disjoint. \\
\textbf{Problem 12.} Let $A$ and $B$ be subsets of a universal set $U$. Prove that if $A \cap B = \emptyset$, then $A \subseteq \overline{B}$.

\subsection{Sets Defined by Number Theory}
\noindent \textbf{Problem 13.} Let $A = \{n \in \mathbb{Z} : n \equiv 1 \pmod 4\}$ and $B = \{n \in \mathbb{Z} : n \equiv 3 \pmod 4\}$. Prove that $A \cap B = \emptyset$. \\
\textbf{Problem 14.} Let $A = \{n \in \mathbb{Z} : 3 \mid n\}$ and $B = \{n \in \mathbb{Z} : 6 \mid n\}$. Prove that $B \subseteq A$. \\
\textbf{Problem 15.} Let $A = \{n \in \mathbb{Z} : n \equiv 2 \pmod 5\}$ and $B = \{n \in \mathbb{Z} : n \equiv 7 \pmod{10}\}$. Prove that $B \subseteq A$. \\
\textbf{Problem 16.} Let $A = \{n \in \mathbb{Z} : 2 \mid n\}$ and $B = \{n \in \mathbb{Z} : 3 \mid n\}$. Prove that $A \cap B = \{n \in \mathbb{Z} : 6 \mid n\}$.

\subsection{Power Sets}
\noindent \textbf{Problem 17.} Let $A$ and $B$ be sets. Prove that $\mathcal{P}(A \cap B) = \mathcal{P}(A) \cap \mathcal{P}(B)$. \\
\textbf{Problem 18.} Let $A$ and $B$ be sets. Prove that $\mathcal{P}(A) \cup \mathcal{P}(B) \subseteq \mathcal{P}(A \cup B)$. \\
\textbf{Problem 19.} Find a rigorous counterexample to show that $\mathcal{P}(A \cup B) = \mathcal{P}(A) \cup \mathcal{P}(B)$ is not generally true. \\
\textbf{Problem 20.} Let $A$ and $B$ be sets. Prove that if $\mathcal{P}(A) \subseteq \mathcal{P}(B)$, then $A \subseteq B$.

\subsection{Generalized Set Operations}
\noindent \textbf{Problem 21.} Let $\{A_i\}_{i \in I}$ be an indexed family of sets within a universal set $U$. Prove the Generalized De Morgan's Law: $\overline{\bigcap_{i \in I} A_i} = \bigcup_{i \in I} \overline{A_i}$. \\
\textbf{Problem 22.} Let $B$ be a set and $\{A_i\}_{i \in I}$ be an indexed family of sets. Prove the Generalized Distributive Law: $B \cap \left(\bigcup_{i \in I} A_i\right) = \bigcup_{i \in I} (B \cap A_i)$.

\subsection{Abstract Modeling: Machine Learning (Symmetric Difference)}
\noindent \textbf{Context:} In statistical learning theory, the symmetric difference $A \Delta B = (A - B) \cup (B - A)$ represents the classification error region between a true positive set $A$ and a predicted positive set $B$. \\
\textbf{Problem 23.} Prove that $A \Delta B = (A \cup B) - (A \cap B)$. \\
\textbf{Problem 24.} Prove that for any set $A$, $A \Delta \emptyset = A$. \\
\textbf{Problem 25.} Prove that for any set $A$, $A \Delta A = \emptyset$. \\
\textbf{Problem 26.} Prove that $A \Delta B = B \Delta A$ (Commutativity of Symmetric Difference).

\subsection{Abstract Modeling: Quantitative Finance}
\noindent \textbf{Context:} Let $\Omega$ be the set of all continuous price paths of a stock over $[0, T]$. Let $\tau_L(\omega)$ be the first time the stock hits a lower bound $L$, and $\tau_H(\omega)$ be the first time it hits an upper bound $H$. Define $A = \{\omega \in \Omega : \tau_L(\omega) \le T\}$ and $B = \{\omega \in \Omega : \tau_H(\omega) \le T\}$. \\
\textbf{Problem 27.} An exotic option pays out only if the stock hits the lower bound but strictly avoids the upper bound before time $T$. Define this success set as $C = A \cap \overline{B}$. Prove rigorously that $C$ and $B$ are disjoint, and that $C \cup (A \cap B) = A$. \\
\textbf{Problem 28.} Let $A_n = \{\omega \in \Omega : \tau_L(\omega) \le T - \frac{1}{n}\}$ for $n \in \mathbb{N}$. Prove that $\bigcup_{n=1}^\infty A_n = \{\omega \in \Omega : \tau_L(\omega) < T\}$.

\subsection{Diagnostic Testing: Identifying Logical Flaws}
\noindent \textbf{Problem 29.} Identify the logical flaw in the following proof and provide the correct logical step. \\
\textit{Claim:} If $A \subseteq B$, then $A \cap C \subseteq B \cap C$. \\
\textit{Flawed Proof:} Assume $A \subseteq B$. We want to show $A \cap C \subseteq B \cap C$. Let $x \in B \cap C$. By definition of intersection, $x \in B$ and $x \in C$. Since $A \subseteq B$, it follows that $x \in A$. Therefore, $x \in A$ and $x \in C$, which means $x \in A \cap C$. Thus, $A \cap C \subseteq B \cap C$. \\
\textbf{Problem 30.} Identify the logical flaw in the following proof and provide the correct logical step. \\
\textit{Claim:} For any set $A$, $\emptyset \cap A = \emptyset$. \\
\textit{Flawed Proof:} We want to show $\emptyset \cap A = \emptyset$. Let $x \in \emptyset \cap A$. By definition of intersection, $x \in \emptyset$ and $x \in A$. Since $x \in \emptyset$, and the empty set is a subset of every set, the intersection must be empty.

\newpage

\section{Detailed Solutions}
\label{sec:solutions}

\subsection{Set Equality via Double Inclusion}

\noindent \textbf{Solution 1.}
\begin{proof}
We use double inclusion to prove $(A \cup B) - (A \cap B) = (A - B) \cup (B - A)$. \\
\textbf{Part 1:} Show $(A \cup B) - (A \cap B) \subseteq (A - B) \cup (B - A)$. \\
Let $x \in (A \cup B) - (A \cap B)$. By definition of set difference, $x \in A \cup B$ and $x \notin A \cap B$. Since $x \in A \cup B$, $x \in A$ or $x \in B$. Since $x \notin A \cap B$, by De Morgan's Laws for logic, $x \notin A$ or $x \notin B$. We proceed by cases based on $x \in A \cup B$:
\begin{itemize}
    \item \textit{Case 1:} $x \in A$. Since $x \notin A$ or $x \notin B$, and we know $x \in A$, it must be that $x \notin B$. Thus, $x \in A$ and $x \notin B$, meaning $x \in A - B$. Therefore, $x \in (A - B) \cup (B - A)$.
    \item \textit{Case 2:} $x \in B$. Since $x \notin A$ or $x \notin B$, and we know $x \in B$, it must be that $x \notin A$. Thus, $x \in B$ and $x \notin A$, meaning $x \in B - A$. Therefore, $x \in (A - B) \cup (B - A)$.
\end{itemize}
In all cases, $x \in (A - B) \cup (B - A)$. \\
\textbf{Part 2:} Show $(A - B) \cup (B - A) \subseteq (A \cup B) - (A \cap B)$. \\
Let $x \in (A - B) \cup (B - A)$. Then $x \in A - B$ or $x \in B - A$.
\begin{itemize}
    \item \textit{Case 1:} $x \in A - B$. Then $x \in A$ and $x \notin B$. Since $x \in A$, $x \in A \cup B$. Since $x \notin B$, $x \notin A \cap B$. Thus, $x \in (A \cup B) - (A \cap B)$.
    \item \textit{Case 2:} $x \in B - A$. Then $x \in B$ and $x \notin A$. Since $x \in B$, $x \in A \cup B$. Since $x \notin A$, $x \notin A \cap B$. Thus, $x \in (A \cup B) - (A \cap B)$.
\end{itemize}
Both inclusions hold, establishing equality.
\end{proof}

\noindent \textbf{Solution 2.}
\begin{proof}
We use double inclusion to prove $A - (B - C) = (A - B) \cup (A \cap C)$. \\
\textbf{Part 1:} Show $A - (B - C) \subseteq (A - B) \cup (A \cap C)$. \\
Let $x \in A - (B - C)$. By definition, $x \in A$ and $x \notin (B - C)$. The statement $x \in B - C$ means $x \in B \land x \notin C$. Negating this via De Morgan's Law yields $x \notin B \lor x \in C$. Since $x \in A$, we can distribute the logical AND: $(x \in A \land x \notin B) \lor (x \in A \land x \in C)$. 
If $x \in A \land x \notin B$, then $x \in A - B$. If $x \in A \land x \in C$, then $x \in A \cap C$. Therefore, $x \in (A - B) \cup (A \cap C)$. \\
\textbf{Part 2:} Show $(A - B) \cup (A \cap C) \subseteq A - (B - C)$. \\
Let $x \in (A - B) \cup (A \cap C)$. Then $x \in A - B$ or $x \in A \cap C$.
\begin{itemize}
    \item \textit{Case 1:} $x \in A - B$. Then $x \in A$ and $x \notin B$. Since $x \notin B$, the statement $x \in B \land x \notin C$ is false, so $x \notin (B - C)$. Thus, $x \in A - (B - C)$.
    \item \textit{Case 2:} $x \in A \cap C$. Then $x \in A$ and $x \in C$. Since $x \in C$, the statement $x \notin C$ is false, so $x \in B \land x \notin C$ is false. Thus $x \notin (B - C)$. Therefore, $x \in A - (B - C)$.
\end{itemize}
Both inclusions hold.
\end{proof}

\noindent \textbf{Solution 3.}
\begin{proof}
We prove $(A - B) - C = A - (B \cup C)$. \\
\textbf{Part 1:} Let $x \in (A - B) - C$. Then $x \in A - B$ and $x \notin C$. This means $(x \in A \land x \notin B) \land x \notin C$. By associativity of logical AND, $x \in A \land (x \notin B \land x \notin C)$. By De Morgan's Laws for logic, $x \notin B \land x \notin C \implies x \notin (B \cup C)$. Thus, $x \in A$ and $x \notin (B \cup C)$, meaning $x \in A - (B \cup C)$. \\
\textbf{Part 2:} Let $x \in A - (B \cup C)$. Then $x \in A$ and $x \notin (B \cup C)$. By De Morgan's Laws, $x \notin (B \cup C) \implies x \notin B \land x \notin C$. Thus, $x \in A \land x \notin B \land x \notin C$. Grouping the first two, $(x \in A \land x \notin B) \implies x \in A - B$. Since $x \notin C$, we have $x \in (A - B) - C$. \\
Both inclusions hold.
\end{proof}

\noindent \textbf{Solution 4.}
\begin{proof}
We prove $A \cap (B - C) = (A \cap B) - (A \cap C)$. \\
\textbf{Part 1:} Let $x \in A \cap (B - C)$. Then $x \in A$ and $x \in B - C$. This means $x \in A \land (x \in B \land x \notin C)$. By associativity and commutativity, $(x \in A \land x \in B) \land x \notin C$. Thus $x \in A \cap B$. Since $x \notin C$, it is impossible for $x \in A \cap C$. Thus $x \notin A \cap C$. Therefore, $x \in (A \cap B) - (A \cap C)$. \\
\textbf{Part 2:} Let $x \in (A \cap B) - (A \cap C)$. Then $x \in A \cap B$ and $x \notin A \cap C$. This means $(x \in A \land x \in B)$ and $(x \notin A \lor x \notin C)$. Since we know $x \in A$ is true, the statement $x \notin A$ is false. For the disjunction $(x \notin A \lor x \notin C)$ to be true, it must be that $x \notin C$. Thus, we have $x \in A \land x \in B \land x \notin C$. Grouping the latter two, $x \in B - C$. Thus $x \in A \cap (B - C)$. \\
Both inclusions hold.
\end{proof}

\subsection{Biconditional Set Properties}

\noindent \textbf{Solution 5.}
\begin{proof}
We prove $A \cup B = B \iff A \subseteq B$. \\
\textbf{Direction 1 ($\implies$):} Assume $A \cup B = B$. We want to show $A \subseteq B$. Let $x \in A$. By definition of union, $x \in A \cup B$. Since $A \cup B = B$, we substitute to get $x \in B$. Thus, $A \subseteq B$. \\
\textbf{Direction 2 ($\impliedby$):} Assume $A \subseteq B$. We want to show $A \cup B = B$. 
First, let $x \in A \cup B$. Then $x \in A$ or $x \in B$. If $x \in A$, since $A \subseteq B$, $x \in B$. If $x \in B$, then trivially $x \in B$. Thus $A \cup B \subseteq B$. 
Second, let $x \in B$. By definition of union, $x \in A \cup B$. Thus $B \subseteq A \cup B$. 
By double inclusion, $A \cup B = B$.
\end{proof}

\noindent \textbf{Solution 6.}
\begin{proof}
We prove $A - B = \emptyset \iff A \subseteq B$. \\
\textbf{Direction 1 ($\implies$):} Assume $A - B = \emptyset$. We want to show $A \subseteq B$. Let $x \in A$. We proceed by contradiction to show $x \in B$. Assume $x \notin B$. Then $x \in A$ and $x \notin B$, which means $x \in A - B$. But $A - B = \emptyset$, so it contains no elements. This is a contradiction. Thus, $x \in B$, meaning $A \subseteq B$. \\
\textbf{Direction 2 ($\impliedby$):} Assume $A \subseteq B$. We want to show $A - B = \emptyset$. We proceed by contradiction. Assume $A - B \neq \emptyset$. Then there exists some $x \in A - B$. By definition, $x \in A$ and $x \notin B$. However, since $A \subseteq B$, $x \in A \implies x \in B$. This contradicts $x \notin B$. Thus, our assumption is false, and $A - B = \emptyset$.
\end{proof}

\noindent \textbf{Solution 7.}
\begin{proof}
We prove $A \cap B = A \cup B \iff A = B$. \\
\textbf{Direction 1 ($\implies$):} Assume $A \cap B = A \cup B$. We want to show $A = B$. 
Let $x \in A$. Then $x \in A \cup B$. Since $A \cup B = A \cap B$, $x \in A \cap B$. By definition of intersection, $x \in B$. Thus $A \subseteq B$. 
Let $x \in B$. Then $x \in A \cup B$. Since $A \cup B = A \cap B$, $x \in A \cap B$. By definition, $x \in A$. Thus $B \subseteq A$. 
Therefore, $A = B$. \\
\textbf{Direction 2 ($\impliedby$):} Assume $A = B$. Then $A \cap B = A \cap A = A$. Also, $A \cup B = A \cup A = A$. Therefore, $A \cap B = A \cup B$.
\end{proof}

\noindent \textbf{Solution 8.}
\begin{proof}
We prove $A \subseteq B \iff \overline{B} \subseteq \overline{A}$. \\
\textbf{Direction 1 ($\implies$):} Assume $A \subseteq B$. Let $x \in \overline{B}$. By definition of complement, $x \notin B$. We proceed by contradiction to show $x \in \overline{A}$. Assume $x \notin \overline{A}$. Then $x \in A$. Since $A \subseteq B$, $x \in A \implies x \in B$. This contradicts $x \notin B$. Thus, $x \in \overline{A}$, meaning $\overline{B} \subseteq \overline{A}$. \\
\textbf{Direction 2 ($\impliedby$):} Assume $\overline{B} \subseteq \overline{A}$. Let $x \in A$. By definition, $x \notin \overline{A}$. We proceed by contradiction to show $x \in B$. Assume $x \notin B$. Then $x \in \overline{B}$. Since $\overline{B} \subseteq \overline{A}$, $x \in \overline{B} \implies x \in \overline{A}$. This contradicts $x \notin \overline{A}$. Thus, $x \in B$, meaning $A \subseteq B$.
\end{proof}

\subsection{Empty Set and Disjointness Proofs}

\noindent \textbf{Solution 9.}
\begin{proof}
Assume $A \cap B \neq \emptyset$ and $A \subseteq C$. Since $A \cap B \neq \emptyset$, there exists an element $x \in A \cap B$. By definition, $x \in A$ and $x \in B$. Since $x \in A$ and $A \subseteq C$, it follows that $x \in C$. Therefore, we have $x \in C$ and $x \in B$, which means $x \in C \cap B$. Since there exists at least one element in $C \cap B$, $C \cap B \neq \emptyset$.
\end{proof}

\noindent \textbf{Solution 10.}
\begin{proof}
Assume $A \cup B = \emptyset$. We want to prove $A = \emptyset$ and $B = \emptyset$. We proceed by contradiction. Assume it is not the case that $A = \emptyset$ and $B = \emptyset$. By De Morgan's Laws for logic, this means $A \neq \emptyset$ or $B \neq \emptyset$. 
\textit{Case 1:} $A \neq \emptyset$. Then there exists $x \in A$. By definition of union, $x \in A \cup B$. This contradicts $A \cup B = \emptyset$. 
\textit{Case 2:} $B \neq \emptyset$. Then there exists $x \in B$. By definition of union, $x \in A \cup B$. This contradicts $A \cup B = \emptyset$. 
In all cases, we reach a contradiction. Thus, $A = \emptyset$ and $B = \emptyset$.
\end{proof}

\noindent \textbf{Solution 11.}
\begin{proof}
We want to prove $A \cap (B - A) = \emptyset$. We proceed by contradiction. Assume $A \cap (B - A) \neq \emptyset$. Then there exists an element $x \in A \cap (B - A)$. By definition of intersection, $x \in A$ and $x \in B - A$. By definition of set difference, $x \in B - A \implies x \in B \land x \notin A$. Thus, we have $x \in A$ and $x \notin A$, which is a logical contradiction. Therefore, our assumption is false, and $A \cap (B - A) = \emptyset$.
\end{proof}

\noindent \textbf{Solution 12.}
\begin{proof}
Assume $A \cap B = \emptyset$. We want to show $A \subseteq \overline{B}$. Let $x \in A$. We must show $x \in \overline{B}$, which means $x \notin B$. We proceed by contradiction. Assume $x \in B$. Since $x \in A$ and $x \in B$, $x \in A \cap B$. However, $A \cap B = \emptyset$, meaning it contains no elements. This is a contradiction. Therefore, $x \notin B$, so $x \in \overline{B}$. Thus, $A \subseteq \overline{B}$.
\end{proof}

\subsection{Sets Defined by Number Theory}

\noindent \textbf{Solution 13.}
\begin{proof}
Let $A = \{n \in \mathbb{Z} : n \equiv 1 \pmod 4\}$ and $B = \{n \in \mathbb{Z} : n \equiv 3 \pmod 4\}$. We proceed by contradiction. Assume $A \cap B \neq \emptyset$. Then there exists $n \in A \cap B$. Thus $n \in A$ and $n \in B$. 
From $n \in A$, $n = 4k + 1$ for some $k \in \mathbb{Z}$. 
From $n \in B$, $n = 4j + 3$ for some $j \in \mathbb{Z}$. 
Equating the two expressions: $4k + 1 = 4j + 3 \implies 4k - 4j = 2 \implies 4(k - j) = 2 \implies 2(k - j) = 1$. 
Since $k, j \in \mathbb{Z}$, $k - j$ is an integer. This equation states that 2 times an integer equals 1, meaning 1 is an even number. This is a contradiction. Thus, $A \cap B = \emptyset$.
\end{proof}

\noindent \textbf{Solution 14.}
\begin{proof}
Let $A = \{n \in \mathbb{Z} : 3 \mid n\}$ and $B = \{n \in \mathbb{Z} : 6 \mid n\}$. We want to show $B \subseteq A$. Let $x \in B$. By definition, $6 \mid x$, which means $x = 6k$ for some integer $k$. We can rewrite this as $x = 3(2k)$. Since $k$ is an integer, $2k$ is an integer. Let $m = 2k$. Then $x = 3m$ for $m \in \mathbb{Z}$. By definition, this means $3 \mid x$, so $x \in A$. Therefore, $B \subseteq A$.
\end{proof}

\noindent \textbf{Solution 15.}
\begin{proof}
Let $A = \{n \in \mathbb{Z} : n \equiv 2 \pmod 5\}$ and $B = \{n \in \mathbb{Z} : n \equiv 7 \pmod{10}\}$. We want to show $B \subseteq A$. Let $x \in B$. Then $x \equiv 7 \pmod{10}$, meaning $x = 10k + 7$ for some $k \in \mathbb{Z}$. We want to show $x \in A$, which requires $x = 5m + 2$ for some $m \in \mathbb{Z}$. 
We manipulate our expression for $x$: $x = 10k + 7 = 10k + 5 + 2 = 5(2k + 1) + 2$. 
Since $k \in \mathbb{Z}$, $2k + 1$ is an integer. Let $m = 2k + 1$. Then $x = 5m + 2$. Thus, $x \equiv 2 \pmod 5$, meaning $x \in A$. Therefore, $B \subseteq A$.
\end{proof}

\noindent \textbf{Solution 16.}
\begin{proof}
Let $A = \{n \in \mathbb{Z} : 2 \mid n\}$, $B = \{n \in \mathbb{Z} : 3 \mid n\}$, and $C = \{n \in \mathbb{Z} : 6 \mid n\}$. We want to show $A \cap B = C$. \\
\textbf{Part 1 ($C \subseteq A \cap B$):} Let $x \in C$. Then $6 \mid x$, so $x = 6k$ for $k \in \mathbb{Z}$. $x = 2(3k)$, so $2 \mid x \implies x \in A$. $x = 3(2k)$, so $3 \mid x \implies x \in B$. Thus $x \in A \cap B$. \\
\textbf{Part 2 ($A \cap B \subseteq C$):} Let $x \in A \cap B$. Then $2 \mid x$ and $3 \mid x$. Since $\gcd(2,3) = 1$, by number theory properties, their product divides $x$. Formally, $x = 2k$. Since $3 \mid x$, $3 \mid 2k$. By Euclid's Lemma, since $\gcd(3,2)=1$, $3 \mid k$. So $k = 3m$ for $m \in \mathbb{Z}$. Substituting back, $x = 2(3m) = 6m$. Thus $6 \mid x$, meaning $x \in C$. \\
Therefore, $A \cap B = C$.
\end{proof}

\subsection{Power Sets}

\noindent \textbf{Solution 17.}
\begin{proof}
We prove $\mathcal{P}(A \cap B) = \mathcal{P}(A) \cap \mathcal{P}(B)$. \\
Let $X \in \mathcal{P}(A \cap B)$. By definition of the power set, $X \subseteq A \cap B$. This means for all $x \in X$, $x \in A \cap B$, so $x \in A$ and $x \in B$. Thus, $X \subseteq A$ and $X \subseteq B$. By definition of the power set, $X \in \mathcal{P}(A)$ and $X \in \mathcal{P}(B)$. Therefore, $X \in \mathcal{P}(A) \cap \mathcal{P}(B)$. \\
The steps are perfectly reversible: If $X \in \mathcal{P}(A) \cap \mathcal{P}(B)$, then $X \in \mathcal{P}(A)$ and $X \in \mathcal{P}(B)$, meaning $X \subseteq A$ and $X \subseteq B$. Thus for all $x \in X$, $x \in A$ and $x \in B$, meaning $x \in A \cap B$. Thus $X \subseteq A \cap B$, so $X \in \mathcal{P}(A \cap B)$. \\
Equality holds.
\end{proof}

\noindent \textbf{Solution 18.}
\begin{proof}
We prove $\mathcal{P}(A) \cup \mathcal{P}(B) \subseteq \mathcal{P}(A \cup B)$. \\
Let $X \in \mathcal{P}(A) \cup \mathcal{P}(B)$. By definition of union, $X \in \mathcal{P}(A)$ or $X \in \mathcal{P}(B)$. 
\textit{Case 1:} $X \in \mathcal{P}(A)$. Then $X \subseteq A$. Since $A \subseteq A \cup B$, by transitivity of subsets, $X \subseteq A \cup B$. Thus $X \in \mathcal{P}(A \cup B)$. 
\textit{Case 2:} $X \in \mathcal{P}(B)$. Then $X \subseteq B$. Since $B \subseteq A \cup B$, by transitivity, $X \subseteq A \cup B$. Thus $X \in \mathcal{P}(A \cup B)$. 
In both cases, $X \in \mathcal{P}(A \cup B)$.
\end{proof}

\noindent \textbf{Solution 19.}
\begin{proof}
We provide a counterexample to $\mathcal{P}(A \cup B) = \mathcal{P}(A) \cup \mathcal{P}(B)$. \\
Let $A = \{1\}$ and $B = \{2\}$. \\
Then $A \cup B = \{1, 2\}$. \\
The power set of the union is $\mathcal{P}(A \cup B) = \{\emptyset, \{1\}, \{2\}, \{1, 2\}\}$. \\
The individual power sets are $\mathcal{P}(A) = \{\emptyset, \{1\}\}$ and $\mathcal{P}(B) = \{\emptyset, \{2\}\}$. \\
Their union is $\mathcal{P}(A) \cup \mathcal{P}(B) = \{\emptyset, \{1\}, \{2\}\}$. \\
Notice that $\{1, 2\} \in \mathcal{P}(A \cup B)$, but $\{1, 2\} \notin \mathcal{P}(A) \cup \mathcal{P}(B)$. Thus, the sets are not equal.
\end{proof}

\noindent \textbf{Solution 20.}
\begin{proof}
Assume $\mathcal{P}(A) \subseteq \mathcal{P}(B)$. We want to show $A \subseteq B$. \\
We know that every set is a subset of itself, so $A \subseteq A$. By definition of the power set, this implies $A \in \mathcal{P}(A)$. \\
Since we assumed $\mathcal{P}(A) \subseteq \mathcal{P}(B)$, and $A \in \mathcal{P}(A)$, it follows that $A \in \mathcal{P}(B)$. \\
By definition of the power set, $A \in \mathcal{P}(B)$ means $A \subseteq B$. This completes the proof.
\end{proof}

\subsection{Generalized Set Operations}

\noindent \textbf{Solution 21.}
\begin{proof}
We prove $\overline{\bigcap_{i \in I} A_i} = \bigcup_{i \in I} \overline{A_i}$. \\
Let $x \in \overline{\bigcap_{i \in I} A_i}$. By definition of complement, $x \notin \bigcap_{i \in I} A_i$. By definition of generalized intersection, it is false that for all $i \in I$, $x \in A_i$. Therefore, there exists at least one $k \in I$ such that $x \notin A_k$. By definition of complement, there exists $k \in I$ such that $x \in \overline{A_k}$. By definition of generalized union, $x \in \bigcup_{i \in I} \overline{A_i}$. \\
The steps are perfectly reversible: If $x \in \bigcup_{i \in I} \overline{A_i}$, there exists $k \in I$ such that $x \in \overline{A_k}$, meaning $x \notin A_k$. Thus it is false that $x$ is in all $A_i$, so $x \notin \bigcap_{i \in I} A_i$, meaning $x \in \overline{\bigcap_{i \in I} A_i}$.
\end{proof}

\noindent \textbf{Solution 22.}
\begin{proof}
We prove $B \cap \left(\bigcup_{i \in I} A_i\right) = \bigcup_{i \in I} (B \cap A_i)$. \\
Let $x \in B \cap \left(\bigcup_{i \in I} A_i\right)$. Then $x \in B$ and $x \in \bigcup_{i \in I} A_i$. By definition of generalized union, there exists some $k \in I$ such that $x \in A_k$. Since $x \in B$ and $x \in A_k$, we have $x \in B \cap A_k$. Since $x$ is in at least one such intersection, $x \in \bigcup_{i \in I} (B \cap A_i)$. \\
Conversely, let $x \in \bigcup_{i \in I} (B \cap A_i)$. Then there exists some $k \in I$ such that $x \in B \cap A_k$. This means $x \in B$ and $x \in A_k$. Since $x \in A_k$, $x \in \bigcup_{i \in I} A_i$. Since $x \in B$ as well, $x \in B \cap \left(\bigcup_{i \in I} A_i\right)$. \\
Equality holds.
\end{proof}

\subsection{Abstract Modeling: Machine Learning}

\noindent \textbf{Solution 23.}
\begin{proof}
We prove $A \Delta B = (A \cup B) - (A \cap B)$. \\
By definition, $A \Delta B = (A - B) \cup (B - A)$. \\
Let $x \in (A - B) \cup (B - A)$. Then $x \in A - B$ or $x \in B - A$. 
If $x \in A - B$, $x \in A$ and $x \notin B$. Since $x \in A$, $x \in A \cup B$. Since $x \notin B$, $x \notin A \cap B$. Thus $x \in (A \cup B) - (A \cap B)$. 
If $x \in B - A$, $x \in B$ and $x \notin A$. Since $x \in B$, $x \in A \cup B$. Since $x \notin A$, $x \notin A \cap B$. Thus $x \in (A \cup B) - (A \cap B)$. \\
Conversely, let $x \in (A \cup B) - (A \cap B)$. Then $x \in A \cup B$ and $x \notin A \cap B$. Since $x \in A \cup B$, $x \in A$ or $x \in B$. Since $x \notin A \cap B$, $x \notin A$ or $x \notin B$. 
If $x \in A$, then we must have $x \notin B$ (otherwise $x \in A \cap B$). Thus $x \in A - B$. 
If $x \in B$, then we must have $x \notin A$. Thus $x \in B - A$. 
In either case, $x \in (A - B) \cup (B - A)$.
\end{proof}

\noindent \textbf{Solution 24.}
\begin{proof}
We prove $A \Delta \emptyset = A$. \\
By definition, $A \Delta \emptyset = (A - \emptyset) \cup (\emptyset - A)$. \\
$A - \emptyset = \{x : x \in A \land x \notin \emptyset\}$. Since $x \notin \emptyset$ is always true, this is just $A$. \\
$\emptyset - A = \{x : x \in \emptyset \land x \notin A\}$. Since $x \in \emptyset$ is always false, this is $\emptyset$. \\
Thus, $A \Delta \emptyset = A \cup \emptyset = A$.
\end{proof}

\noindent \textbf{Solution 25.}
\begin{proof}
We prove $A \Delta A = \emptyset$. \\
By definition, $A \Delta A = (A - A) \cup (A - A)$. \\
$A - A = \{x : x \in A \land x \notin A\}$. This is a logical contradiction, so $A - A = \emptyset$. \\
Thus, $A \Delta A = \emptyset \cup \emptyset = \emptyset$.
\end{proof}

\noindent \textbf{Solution 26.}
\begin{proof}
We prove $A \Delta B = B \Delta A$. \\
By definition, $A \Delta B = (A - B) \cup (B - A)$. \\
Since the set union operation is commutative ($X \cup Y = Y \cup X$), we can swap the terms: \\
$(A - B) \cup (B - A) = (B - A) \cup (A - B)$. \\
By definition, $(B - A) \cup (A - B) = B \Delta A$. \\
Thus, $A \Delta B = B \Delta A$.
\end{proof}

\subsection{Abstract Modeling: Quantitative Finance}

\noindent \textbf{Solution 27.}
\begin{proof}
We are given $C = A \cap \overline{B}$. \\
\textbf{Part 1 (Disjointness):} We must show $C \cap B = \emptyset$. \\
$C \cap B = (A \cap \overline{B}) \cap B$. By associativity of intersection, this is $A \cap (\overline{B} \cap B)$. Since a set and its complement are disjoint, $\overline{B} \cap B = \emptyset$. Thus, $A \cap \emptyset = \emptyset$. Therefore, $C$ and $B$ are disjoint. \\
\textbf{Part 2 (Union):} We must show $C \cup (A \cap B) = A$. \\
$C \cup (A \cap B) = (A \cap \overline{B}) \cup (A \cap B)$. By the Distributive Law of sets, this equals $A \cap (\overline{B} \cup B)$. Since the union of a set and its complement is the universal space $\Omega$, we have $A \cap \Omega$. Since $A \subseteq \Omega$, $A \cap \Omega = A$.
\end{proof}

\noindent \textbf{Solution 28.}
\begin{proof}
Let $A_n = \{\omega \in \Omega : \tau_L(\omega) \le T - \frac{1}{n}\}$. We want to prove $\bigcup_{n=1}^\infty A_n = \{\omega \in \Omega : \tau_L(\omega) < T\}$. Let $S = \{\omega \in \Omega : \tau_L(\omega) < T\}$. \\
\textbf{Part 1:} Let $\omega \in \bigcup_{n=1}^\infty A_n$. By definition of union, there exists some $k \in \mathbb{N}$ such that $\omega \in A_k$. This means $\tau_L(\omega) \le T - \frac{1}{k}$. Since $k \ge 1$, $\frac{1}{k} > 0$, which implies $T - \frac{1}{k} < T$. Therefore, $\tau_L(\omega) < T$, meaning $\omega \in S$. \\
\textbf{Part 2:} Let $\omega \in S$. Then $\tau_L(\omega) < T$. Let $\epsilon = T - \tau_L(\omega)$. Since $\tau_L(\omega) < T$, $\epsilon > 0$. By the Archimedean property of real numbers, there exists an integer $k \in \mathbb{N}$ such that $\frac{1}{k} \le \epsilon$. Thus, $\frac{1}{k} \le T - \tau_L(\omega)$, which rearranges to $\tau_L(\omega) \le T - \frac{1}{k}$. This means $\omega \in A_k$. Since $\omega$ is in at least one $A_k$, $\omega \in \bigcup_{n=1}^\infty A_n$. \\
Both inclusions hold.
\end{proof}

\subsection{Diagnostic Testing: Identifying Logical Flaws}

\noindent \textbf{Solution 29.}
\begin{proof}
\textbf{Flaw Identification:} The proof commits the \textit{Converse Error} (or Affirming the Consequent). To prove $A \cap C \subseteq B \cap C$, the proof MUST start by assuming an arbitrary element is in the left-hand set: "Let $x \in A \cap C$." Instead, the flawed proof starts with "Let $x \in B \cap C$" and attempts to work backward, which proves $B \cap C \subseteq A \cap C$ (if the steps were valid, which they aren't, because $x \in B$ does not imply $x \in A$). \\
\textbf{Correct Logical Step:} \\
Assume $A \subseteq B$. We want to show $A \cap C \subseteq B \cap C$. \\
Let $x \in A \cap C$. By definition of intersection, $x \in A$ and $x \in C$. \\
Since $x \in A$ and $A \subseteq B$, it follows that $x \in B$. \\
Now we have $x \in B$ and $x \in C$. By definition of intersection, $x \in B \cap C$. \\
Therefore, $A \cap C \subseteq B \cap C$.
\end{proof}

\noindent \textbf{Solution 30.}
\begin{proof}
\textbf{Flaw Identification:} The proof falls into the \textit{Empty Set Trap}. It begins with "Let $x \in \emptyset \cap A$." This assumes that the intersection is non-empty and that such an element $x$ can be instantiated. Since the goal is to prove the set is empty, instantiating an element from it is logically invalid. \\
\textbf{Correct Logical Step:} \\
We want to show $\emptyset \cap A = \emptyset$. We proceed by contradiction. \\
Assume $\emptyset \cap A \neq \emptyset$. Then there exists some element $x \in \emptyset \cap A$. \\
By definition of intersection, $x \in \emptyset$ and $x \in A$. \\
However, the statement $x \in \emptyset$ is a logical contradiction, as the empty set contains no elements by definition. \\
Therefore, our assumption must be false, and $\emptyset \cap A = \emptyset$.
\end{proof}

\newpage

\section{Concept Checklist \& Question Mapping (Repeated)}
\label{sec:mapping_end}
\begin{itemize}
    \item \textbf{C1: Set Equality via Double Inclusion} (Problems 1--4)
    \item \textbf{C2: Biconditional Set Properties} (Problems 5--8)
    \item \textbf{C3: Empty Set, Disjointness, and Contrapositive/Contradiction} (Problems 9--12)
    \item \textbf{C4: Sets Defined by Number Theory (Modular Arithmetic \& Divisibility)} (Problems 13--16)
    \item \textbf{C5: Power Sets ($\mathcal{P}(A)$) and Element vs. Subset Logic} (Problems 17--20)
    \item \textbf{C6: Generalized Set Operations (Infinite Unions/Intersections)} (Problems 21--22)
    \item \textbf{C7: Abstract Modeling: Machine Learning (Symmetric Difference)} (Problems 23--26)
    \item \textbf{C8: Abstract Modeling: Quantitative Finance (Probability Spaces \& Hitting Times)} (Problems 27--28)
    \item \textbf{C9: Diagnostic Testing: Identifying Logical Flaws} (Problems 29--30)
\end{itemize}

\end{document}